\subsection{收缩定理的证明}
\label{app:tensor}

在定理~\ref{thm:tensor-contraction} 的整个证明中，$\log$ 表示以二为底的对数，$\ln$ 表示自然
对数。

\subsubsection{系数环与精确算术}
\label{subsec:tensor-arithmetic}

记 $\lambda=\lceil\log_2 s\rceil$，$e=4\lambda$。
选取首一多项式 $f\in\mathbb Z[U]$，次数为 $e$，其首项系数以下的系数
属于 $\{0,1\}$，且其模二约化不可约。这可以在关于 $s$ 的多项式时间内
确定性地完成：枚举 $2^e$ 个候选，并对每个次数不超过 $e/2$ 的正次数
首一多项式检验 $\mathbb F_2$ 上的整除性。有限域在每个正次数上
都存在不可约多项式~\cite{vzGG13}。
令
\[
 \mathcal O=\mathbb Z_{(2)}[U]/(f),\qquad
 \mathcal K=\mathbb Q[U]/(f),\qquad
 \mathcal A=\mathcal O[V]/(V^\ell),\qquad
 \mathcal A_p=\mathcal A/(2^p).
\]
这里 $\mathbb Z_{(2)}$ 由分母为奇数的有理数组成。模二不可约蕴含
在 $\mathbb Q$ 上不可约，因此 $\mathcal K$ 是域。对 $\mathcal K$ 中的非零
元素，令 $v_2$ 为其在 $1,U,\ldots,U^{e-1}$ 下坐标的通常二进赋值的
最小值。$\mathcal O$ 模二的约化是域 $\mathbb F_{2^e}$。
因此约去最小的坐标赋值即表明
\[
 v_2(xy)=v_2(x)+v_2(y)\qquad(x,y\ne0).
\]
$\mathcal O$ 中任何模二非零的元素都是单位，因为其乘法矩阵的行列式
在 $\mathbb Z_{(2)}$ 中可逆，故逆矩阵的元素属于该环。

$\mathcal A_p$ 中的算术在 $s,p,\ell$ 上花费多项式时间与空间。我们使用
基 $1,U,\ldots,U^{e-1}$ 与 $1,V,\ldots,V^{\ell-1}$，并把每个坐标表示为
零到 $2^p-1$ 之间的整数剩余。标量单位的求逆可以通过模二求逆再作
Newton 提升完成，也可以通过其乘法矩阵完成，所用的是精确的多项式
与整数运算~\cite{vzGG13}。$\mathcal A$ 中元素的可除性与精度始终是
逐坐标的。

$2^e$ 个剩余域元素的二进制坐标提升给出两两差为单位的内插点。
因此在 $\mathcal A_p$ 上，在任意所需数量的这类点上作 Lagrange
插值是合法的。我们将使用少于 $s^4$ 个点。
初始元素属于 $\mathcal B_{P_0,\ell}$ 的收缩在 $\mathcal A_{P_0}$ 中计算后
仍属于该子环，因为它是这些元素的多项式。

当精度提高时，我们逐个提升每个张量元素。对于最初模 $2^p$ 指定的
元素，其\emph{典范提升}是由其所有坐标的整数代表元得到的 $\mathcal A$
中元素。典范提升不必保持加法或乘法。每次递归调用都按其调用的
精度收缩指定的张量。为定义其形变，它首先作这种逐元素的提升
选择。该调用中随后的所有求值都指向这些相同的提升。
更明确地说，每个张量描述都标记有其创建精度 $p_c$（原始张量为 $P_0$）。
在之后每个精度下，其元素都模 $2^{p_c}$ 求值，并使用相同的整数
代表元。精度沿递归路径严格递增，故这些代表元在之后每个精度下
也是典范代表元。特别地，
包张量在整个递归过程中保持其创建时的代表元。

\subsubsection{参数与碰撞打包}
\label{subsec:tensor-packing}

以下所有估计对充分大的 $s$ 成立，其阈值只依赖于固定常数。较小的
有限尺寸范围由直接求和处理。参数不等式可以直接检验。令
\[
 H=\left\lceil\frac{s}{\lfloor\sqrt\lambda\rfloor}\right\rceil,
 \qquad r=64,\qquad \varepsilon=\frac1{128\cdot1024},
 \qquad \Delta=\lceil s^{9/10}\rceil.
\]
边数 $M\le H$ 的调用是叶，直接求值。在深度 $d$ 的活动调用处（根在
深度零），使用
\begin{equation}\label{eq:tensor-parameters}
 u=u_d=\left\lfloor\frac{\varepsilon\lambda}{2^d}\right\rfloor,
 \qquad a=a_d=2\sum_{j=0}^d u_j,\qquad
 a_{-1}=0,\qquad L=2^a,\qquad D_{\rm arr}=r^{2a}.
\end{equation}
我们将安排每个子节点至多有 $M/8$ 条边且最大度数至多 $\Delta$。
这些断言表明，在每条活动路径上，
\[
 H<M\le Cs\,8^{-d},\qquad
 d\le\frac{\log_2\lambda}{6}+O_C(1).
\]
因此
\begin{equation}\label{eq:tensor-parameter-bounds}
 u=\Theta_C(\lambda/2^d),\quad
 \lambda^{4/5}\le u\le M,\quad
 a\le4\varepsilon\lambda,\quad
 L,D_{\rm arr}\le s^{1/1000}
\end{equation}
对所有充分大的 $s$ 成立。特别地，在每个活动调用处
\begin{equation}\label{eq:tensor-packing-conditions}
 M\ge2D_{\rm arr}^{\,2}\Delta,\qquad
 M\ge L D_{\rm arr},\qquad 2L\le\Delta.
\end{equation}
例如，第一个不等式右端为
$O(s^{9/10+2/1000})=o(H)$。其他不等式以同样方式得出。

\begin{lemma}[打包成包]\label{lem:tensor-packets}
在 \eqref{eq:tensor-packing-conditions} 下，用 $[r]^a$ 的元素标记网络每侧的
顶点。存在一个确定性可构造的标记，以及一个把边嵌入到含 $L$ 个槽的
包中的嵌入，具有以下性质。
\begin{enumerate}
\item 至多有 $3M/L$ 个包。每条边在其行端与列端占据相同的槽。
      未使用的槽没有端口。
\item 在每个包和每个坐标中，行符号与列符号的无序对是固定的。
      当这两个符号不同时，槽号的第 $i$ 位指定它们在坐标 $i$ 中的
      取向，其中 $0\le i<a$。
\item 对每个坐标和每个符号，每侧至多有 $3M/r$ 个实端口属于具有
      该符号的顶点。
\end{enumerate}
该构造在 $s+\mathcal L$ 上使用多项式时间与空间。
\end{lemma}
\begin{proof}
对每对标记，令 $X_\gamma$ 为具有该标记对的边数。独立的均匀顶点
标记给出
\[
 \mathbb E\sum_\gamma (X_\gamma-M/D_{\rm arr})^2\le2M\Delta.
\]
事实上，展开把该期望写成有序边对上其标记一致概率减去
$1/D_{\rm arr}$ 的求和。端点不交的边对贡献为零。
至多有 $2M\Delta$ 个有序边对共享一个端点（包括相同的边），
每对至多贡献一。
条件期望给出具有该界的标记。为实现它，对下一个顶点的 $r^a$ 个
标记逐一尝试，并通过枚举边对计算条件期望。每个概率至多涉及
四个顶点，故这是精确的多项式时间计算。这是通常的
条件期望方法~\cite{Raghavan88}。

每个负载至多为 $M/D_{\rm arr}+\sqrt{2M\Delta}$。
把所有标记对类填充到共同的槽数 $H_{\rm pad}$，即最大负载向上取整
到 $L$ 的倍数。假设蕴含
\[
 H_{\rm pad}\le M/D_{\rm arr}+\sqrt{2M\Delta}+L
              \le3M/D_{\rm arr}.
\]
固定每个坐标中符号的无序对。若有 $a'$ 个坐标的符号不相等，
其取向给出 $2^{a'}$ 个标记对。用每个包中来自这些标记的
$2^{a-a'}$ 个槽把它们的填充槽划分为包。
这可行是因为 $H_{\rm pad}$ 可被 $L$ 整除。
把一个包的位置从 $0$ 编号到 $L-1$。其第 $i$ 位在第 $i$ 对不相等时
指定取向，在该对相等时指定副本。在每个包中始终使用一种取向约定。

有 $H_{\rm pad}D_{\rm arr}/L\le3M/L$ 个包。在固定坐标中，给定符号
恰出现在每侧所有填充标记对的 $1/r$ 中。因此其实端口数至多为
$H_{\rm pad}D_{\rm arr}/r\le3M/r$。
\end{proof}

\subsubsection{一个 Cauchy 形变与整局部坐标}
\label{subsec:tensor-deformation}

固定一个活动调用、它的包以及它逐元素提升的张量。使用
\[
 S=\{-L,-L+1,\ldots,L-1\},\qquad
 \mathcal C(T)=\prod_{t\in S}(T-t),\qquad \alpha=2^{a+1}.
\]
在一个包内，对其行槽 $x$ 与列槽 $y$，令
\begin{equation}\label{eq:tensor-cauchy}
 W_{xy}(T)=\frac{\mathcal C(T)}
                 {\mathcal C'(0)(T-(y-x))}.
\end{equation}
分母中的因子与分子中的一个因子相消，因为 $y-x\in S$。
因此这是 $\mathbb Q$ 上次数为 $2L-1$ 的多项式。把不同包之间的
元素置为零，并限制到实端口。在原来的边识别下，
$W(0)=I$，因为对角元素为一，而其他每个元素都有一个未相消的
零因子。

令 $A(Y)$ 与 $B(X)$ 分别为提升后的行张量与列张量在独立端口赋值
$Y,X\in\{0,1\}^M$ 上的乘积。记 $l_0(z)=1-z$，$l_1(z)=z$。定义
\begin{equation}\label{eq:tensor-deformed-contraction}
 F(T)=\sum_{X,Y\in\{0,1\}^M} A(Y)B(X)
                 \prod_{\text{行端口 }x}l_{Y_x}((W(T)X)_x).
\end{equation}
其系数属于 $\mathcal K[V]/(V^\ell)$，次数至多为 $M(2L-1)$，
且 $F(0)$ 是提升后张量的收缩。
特别地，把 $F(0)$ 模该调用的精度约化即给出所需答案。

\begin{lemma}\label{lem:tensor-integrality}
对每个 $t\in S$，$W(t+\alpha z)$ 的元素属于 $\mathcal O[z]$，且
$F(t+\alpha z)$ 属于 $\mathcal A[z]$。
给定 $\zeta$ 模 $2^p$ 的 $e$ 个坐标，每个元素
$W_{xy}(t+\alpha\zeta)\bmod 2^p$ 可用 $s,p$ 上多项式多个位运算算出。
\end{lemma}
\begin{proof}
网格 $S$ 是模 $2^{a+1}$ 的完全剩余系。
对 $S$ 中 $t\ne q$，$v_2(t-q)\le a$，且各赋值上非零差的计数
与 $t$ 无关。因此 $v_2(\mathcal C'(t))=v_2(\mathcal C'(0))$。
在
\[
 \frac{\mathcal C(t+\alpha z)}{\mathcal C'(0)}
  =\frac{\alpha z\prod_{q\in S\setminus\{t\}}
                        (t-q+\alpha z)}{\mathcal C'(0)},
\]
中，每个由 $q\ne t$ 索引的因子都按系数被 $2^{v_2(t-q)}$ 整除。
把这些幂与分母相消后，所示多项式被 $2^{a+1}$ 整除。
为得到 \eqref{eq:tensor-cauchy}，只略去由 $y-x$ 索引的那个因子。
若 $y-x=t$，这去掉 $\alpha z$，留下一个整多项式。
否则 $\alpha z$ 保留，而被略去的因子有二进容度 $2^{v_2(t-y+x)}$，
其中 $v_2(t-y+x)\le a$，故结果仍被二整除。代入
\eqref{eq:tensor-deformed-contraction} 即证得关于 $F$ 的断言。

对求值，先用已知整数 $2^{v_2(t-q)}$ 除每个因子
$t-q+\alpha z$（$q\ne t$）。所得线性多项式的两个系数都是整数。
如上把剩余的显式二的幂与分母相消，并将其奇部模 $2^p$ 求逆。
现在我们可以在给定的 $\zeta$ 剩余处直接计算乘积，包括
$\zeta=0$。输入有 $ep$ 位，乘积有 $O(L)$ 个因子。
所有坐标运算在 $s,p$ 上花费多项式时间。
\end{proof}

\subsubsection{渐缩 Hermite 恢复}
\label{subsec:tensor-hermite}

下一个引理分离出精度估计。其陈述中所有幂级数截断都是 $z$ 中的
通常截断。

\begin{lemma}[渐缩 Hermite 恢复]\label{lem:tapered-hermite}
令整数 $u,M$ 满足 $1\le u\le a$ 与 $M\ge u$，并令
$h=\lceil3M/u\rceil$。
对 $t\in S\setminus\{0\}$ 置
\[
 i=v_2(t),\qquad q_t=q(i)=h\max(0,u-i),\qquad D_t=L/2^i.
\]
假设 $F\in(\mathcal K[V]/(V^\ell))[T]$ 的次数至多为
$M(2L-1)$，$F(0)\in\mathcal A$，且
对每个 $t\in S$ 有 $F(t+\alpha z)\in\mathcal A[z]$。
对 $p\ge1$，其值 $F(0)$ 模 $2^p$ 可由
\[
 F(t+\alpha z)\pmod{2^{p_t}},
 \qquad
 p_t=p+8hD_t+8M(a+1),\qquad q_t>0.
\]
恢复，其中只用到每个多项式的前 $q_t$ 个系数。
恢复在 $M,L,\ell,e$ 与各精度上花费多项式时间与空间。
\end{lemma}
\begin{proof}
$S$ 中有 $L$ 个奇点，故
$\sum_{t\ne0}q_t\ge Lhu\ge3ML>M(2L-1)$。又
$h\le4M/u$，故 $q_t\le4M$。
对每个重数为正的点 $t$ 置
\[
 E_t(T)=\prod_{q\in S\setminus\{0,t\}}
                   \left(\frac{T-q}{t-q}\right)^{q_q}.
\]
Hermite 插值给出
\begin{equation}\label{eq:tensor-hermite-origin}
 F(0)=\sum_{t:q_t>0}E_t(0)
       \left(\left[\frac{F(t+\alpha z)}
                            {E_t(t+\alpha z)}\right]_{<q_t}\right)
                            (-t/\alpha).
\end{equation}
为完整起见，把右端的零换为未定元 $T$，把 $-t/\alpha$ 换为
$(T-t)/\alpha$。每个被加项的次数小于 $\sum q_t$。在其自身节点处
它具有规定的射流，而在其他每个节点处它按所需阶数消失。
因此与 $F$ 的差被
$\prod_{t\ne0}(T-t)^{q_t}$ 整除且次数更小，故为零。
这是 Hermite 插值的通常中国剩余构造 \cite[第~5.6~节]{vzGG13}。

$E_t(t+\alpha z)$ 的每个因子都是 $1+\alpha z/(t-q)$ 的幂。
它是整的且常数项为一。
因此其截断到 $q_t$ 的倒数也是整的，故在
\eqref{eq:tensor-hermite-origin} 中作商不损失精度。

固定 $i=v_2(t)$ 与 $D=D_t$。直接由乘积得
\begin{align}
 v_2(E_t(0))
 &=\sum_{q\ne0,t}q_q\bigl(v_2(q)-v_2(t-q)\bigr)\notag\\
 &=D\sum_{j=1}^{a-i}\frac{j}{2^j}
                          \bigl(q(i+j)-q(i)\bigr)
 \ \ge -6Dh. \label{eq:tensor-hermite-loss}
\end{align}
为说明第二个等式，赋值小于 $i$ 的点贡献为零。对 $1\le j\le a-i$，
有 $D/2^j$ 个点满足 $v_2(q)=i+j$，同样多个点满足
$v_2(t-q)=i+j$。后者满足 $v_2(q)=i$，且这些层穷尽了具有该赋值的
除 $t$ 以外的点。这些计数来自完全剩余网格，其非零整数差的
赋值至多为 $a$。最后，
$q(i+j)-q(i)\ge-hj$ 且 $\sum_{j\ge1}j^2/2^j=6$。

次数低于 $q_t$ 的多项式在 $-t/\alpha$ 处求值至多再损失
$4M(a+1)$ 个二的幂。因此把每个商系数替换为其模
$2^{p_t}$ 剩余的任何提升，会使其在
\eqref{eq:tensor-hermite-origin} 中的贡献改变 $2^p\mathcal A$ 中的
一个元素。把这些误差求和不再损失精度。重构的和是整的，因为它与
整值 $F(0)$ 相差 $2^p\mathcal A$ 中的一个元素。

可以在指定基中用有理算术完成最后的组合。乘子
$E_t(0)(-t/\alpha)^j$（$j<q_t$）是多项式多个 $O(a)$ 位有理因子的
乘积（计重数）。因此其分子与分母长度在 $M,L,a$ 中为多项式。
这样的系数只有多项式多个，故即使取其分母的乘积也构成合适的
多项式位公共分母。相加后，约去其二的幂，并把所得整坐标模 $2^p$
约化。这就证明了算法断言。
\end{proof}

在输入精度为 $p$ 的调用处，用 \eqref{eq:tensor-parameters} 中的参数
应用此引理。对每个活动的 $t$，
在 $M(2L-1)+1$ 个不同的剩余域元素的二进制坐标提升处求值并作 Lagrange
插值，得到整个多项式 $F(t+\alpha z)$ 模 $2^{p_t}$。
引理~\ref{lem:tensor-integrality} 允许在 $\mathcal A_{p_t}$ 上作此
插值。由 \eqref{eq:tensor-parameter-bounds}，点数少于 $s^4$。
下面的求值算法把精度 $p_t$ 传给其子节点。

输入精度与子精度满足
\begin{equation}\label{eq:tensor-precision-bound}
 p\le s\lambda^{O(1)}2^{a_{d-1}},
 \qquad
 p_t=O_C(MD_t/u).
\end{equation}
第一个断言在根处成立，因为 $P_0\le Cs$。沿一条路径，精度的增量为
$8hD_t+8M(a+1)$，其中 $h\le4M/u$、$M\le Cs$、$D_t\le2^{a_d}$。
把它们在 $O(\log\lambda)$ 个深度上求和，即证得对每个调用（包括叶）
的第一个断言。对第二个断言，$q_t>0$ 蕴含 $i<u$，从而
\[
 D_t\ge2^{a_d-u}=2^{a_{d-1}+u}.
\]
由于 $M>H$ 且 $u\ge\lambda^{4/5}$，这个额外因子 $2^u$ 支配 $\lambda$
的每个固定幂。因此关于 $p$ 的第一个界以及 $8M(a+1)$ 都是
$O_C(MD_t/u)$。剩余的增量 $8hD_t$ 直接具有该界。
所有精度在 $s$ 中保持多项式，因为
$a_d\le4\varepsilon\lambda$。

渐缩解释了此估计中的因子 $1/u$。
作为比较，给 $u$ 以下的每个层级指定相同的重数 $q_*=\Theta(M)$，
其余层级指定为零。在最后保留的层级 $i=u-1$ 处，恒等式
\eqref{eq:tensor-hermite-loss} 会给出损失
\[
 -q_*D_t\sum_{j=1}^{a-i}\frac{j}{2^j}=-\Theta(MD_t).
\]
而我们的重数每层只改变 $h=\Theta(M/u)$，
给出损失 $O(MD_t/u)$。$a_d$ 的累积选择提供了吸收继承精度与
求值损失所需的额外因子 $2^u$。因此在每个保留层级与每个递归
深度处，较小的界都成立。

\subsubsection{群之间相互作用的模式}
\label{subsec:tensor-modes}

固定一个活动的 $t$ 与 $z$ 的一个插值 $\zeta\in\mathcal O$。在本小节中
$T=t+\alpha\zeta$，所有计算都模 $2^{p_t}$。置 $i=v_2(t)<u$ 与
$D=L/2^i$。
按顶点在坐标 $i$ 中的符号把两侧的顶点分组。
$W(T)$ 中两个端口属于同一群的元素称为局部的，并保持不变。

这种坐标选择使 $W(t)$ 成为局部的。事实上，
除非 $y-x=t$，$W_{xy}(t)=0$。加上一个赋值为 $i$ 的整数会翻转第 $i$ 位，
故存活的边对具有不同的第 $i$ 位。当包在该坐标中的两个符号不同时，
这意味着行符号与列符号一致。相等的符号对
本来就位于一个群中。从 $t$ 移到 $t+\alpha\zeta$ 会引入
我们现在展开的交叉项。

在坐标 $i$ 中两个符号不同的包里，行位置 $x$ 与列位置 $y$ 位于
不同群中，恰好当它们的第 $i$ 位相等。把列划分为瓦片
\[
 y=y_0+2^{i+1}j,\qquad 0\le y_0<2^{i+1}.
\]
列群在一个瓦片上为常数。对一条交叉到该瓦片的行，
\[
 v_2(t-y_0+x)\le i.
\]
否则 $y_0$ 会与 $x+2^i$ 模 $2^{i+1}$ 一致，与其第 $i$ 位相等矛盾。
加上 $\alpha\zeta$ 不改变该赋值。因此几何展开给出
\begin{equation}\label{eq:tensor-mode-expansion}
 W_{xy}(T)=
 \sum_{\substack{b\ge0\\b+1<p_t}}
 \frac{\mathcal C(T)}
      {\mathcal C'(0)(T-y_0+x)}
 \left(\frac{2^{i+1}}{T-y_0+x}\right)^b j^b
 \pmod {2^{p_t}}.
\end{equation}
事实上第一个分式的赋值至少为 $a+1-i\ge1$，而该比在取 $b$ 次幂
之前的赋值至少为一。几何余项被 $2^{p_t}$ 整除。这些表达式可以
通过分解已知的二的幂并只用单位逆来计算。
对于每个这样的包、瓦片和 $b$，引入一个权为 $w_Z=b+1$ 的模式变量 $Z$，并以
\[
 L_Z(X)=\sum_{\text{瓦片中的实列 }y}j^bX_y.
\]
替换它。按惯例 $j^0=1$，包括 $j=0$ 的情形。其源是该瓦片的列组，其目标是与该包坐标对配对的另一个组。公式
\eqref{eq:tensor-mode-expansion} 将每个行形式表示为其局部部分与其目标模式的线性组合之和。
$Z$ 的系数可被 $2^{w_Z}$ 整除。任一权重的变量至多有
\begin{equation}\label{eq:tensor-mode-multiplicity}
 (3M/L)2^{i+1}=6M/D
\end{equation}
个。

在仿射因子
$l_{Y_x}((W(T)X)_x)$ 的乘积中，每个模式单项式
$Z^\beta$ 的系数都可被 $2^{\sum_Zw_Z\beta_Z}$ 整除。
因此只需考虑满足
\begin{equation}\label{eq:tensor-mode-budget}
 \sum_Zw_Z\beta_Z<p_t
\end{equation}
的指数元组。这些元组的数目为
\begin{equation}\label{eq:tensor-partition-count}
 \exp\bigl(O(\sqrt{Mp_t/D})\bigr)
       =\exp\bigl(O_C(M/\sqrt u)\bigr).
\end{equation}
以下是完整的计数估计。对任意 $\theta>0$，
生成乘积与 \eqref{eq:tensor-mode-multiplicity} 将该数目界为
\[
 \exp\left(\theta p_t+\frac{6M}{D}
                   \sum_{j\ge1}-\ln(1-e^{-\theta j})\right).
\]
展开对数并对几何级数求和得
\[
 \sum_{j\ge1}-\ln(1-e^{-\theta j})
  =\sum_{h\ge1}\frac1{h(e^{\theta h}-1)}
  \le\frac1\theta\sum_{h\ge1}\frac1{h^2}.
\]
取 $\theta=\sqrt{M/(Dp_t)}$ 并使用
\eqref{eq:tensor-precision-bound}。所有模式变量及其
指数都有多项式规模的描述。满足
\eqref{eq:tensor-mode-budget} 的元组可按变量排序，并在剩余预算的约束下依次选取
每个指数，从而以多项式延迟和空间枚举。每个可行的部分
选择都有全零的延续，因此访问的部分选择数至多是完整元组数
乘以变量数，至多相差一个加性常数。

\subsubsection{子网络及其描述}
\label{subsec:tensor-children}

对于 \eqref{eq:tensor-mode-budget} 中的一个元组 $\beta$，为每个组 $g\in[r]$ 构造
一个子网络。把当前网络中来自两侧的所有当前顶点都放在
子网络的一侧。保留它们的端口及其选定的提升项，此时已对 $2^{p_t}$ 取模。
在另一侧，为每个包放置一个顶点，并向该包每个保留下来的端口
连一条边。其指标恰为该包中组 $g$ 的行位 $Y$ 与列位 $X$。

为定义这个新的包张量，取包中组 $g$ 的行因子 $l_{Y_x}$ 的乘积，其中各因子
的自变量用局部列位与形式模式表示。对所有以 $g$ 为目标的模式，提取指数为
$\beta_Z$ 的系数。将其乘以
\[
 \prod_{\substack{Z\text{ 在该包中}\\
                  \text{source}(Z)=g}}L_Z(X)^{\beta_Z}.
\]
空乘积等于一。若没有目标行，则只要某个目标指数为正，提取出的
系数就是零。若 $Z$ 的源瓦片没有实列，则 $L_Z(X)=0$，其
幂在 $\beta_Z=0$ 时贡献一，否则贡献零。

每个模式恰有一个目标和一个源。因此系数提取可沿包-组行乘积分解，
且每个模式的代换只依赖于其源组中的位。将
\eqref{eq:tensor-deformed-contraction} 中的乘积展开，再对组内相互独立的位求和，
即证明精确恒等式
\begin{equation}\label{eq:tensor-recursion-identity}
 F(t+\alpha\zeta)
   =\sum_{\beta:\,\sum w_Z\beta_Z<p_t}
                       \prod_{g=1}^r Z_{g,\beta}
                  \pmod {2^{p_t}},
\end{equation}
其中 $Z_{g,\beta}$ 是刚才定义的子网络的收缩。被省略的单项式由其系数整除性而消失。

由引理~\ref{lem:tensor-packets}，一个子网络至多有
$6M/r=3M/32<M/8$ 条边。旧顶点的度不变，而
新的包顶点至多有 $2L\le\Delta$ 个端口。因此
\eqref{eq:tensor-parameter-bounds} 中使用的不变量归纳成立。

所有新张量都是隐式给出的。要计算某个包张量的一项，展开其至多 $L$ 个仿射行
因子的乘积。每个因子关于 $s$ 的项数都是多项式的，包括其局部项。
枚举这些选择、检查所需指数并对所得各项求和，需要
\[
 \exp(O(L\log s))
\]
时间和多项式空间。因子 $L_Z(X)^{\beta_Z}$ 通过在 $\mathcal A_{p_t}$ 中通常的反复平方
计算。由于 $L\le s^{1/1000}$，这一项的计算代价为
$\exp(O(s^{9/10}))$。

新张量的描述由其包数据、元组 $\beta$、插值取值和创建精度
组成。这些数据连同局部核足以直接计算其项。保留下来的张量
按其原始描述计算；若它是在递归中更早创建的，则按其自身的
创建时描述计算。其在该精度下的剩余类随后被规范地
提升到当前精度。旧描述是不可变的，并通过引用共享。因此项计算
例程的任何复合都不会随深度累积，且所有项计算例程
仍具有所需的时间与空间界。
沿一条路径只添加多项式多个顶点和多项式多的新
描述数据。

\subsubsection{时间与空间界}
\label{subsec:tensor-cost}

\begin{proof}[定理~\ref{thm:tensor-contraction} 的证明]
将原始项嵌入 $\mathcal A_{P_0}$ 并执行上述递归算法。在一个活动调用处，对每个
重数为正的节点和每个所需的剩余类互异插值取值，通过其子收缩计算 \eqref{eq:tensor-recursion-identity}。
对 $F(t+\alpha z)$ 插值并应用
引理~\ref{lem:tapered-hermite}。在叶节点处，直接对所有
布尔边赋值求和。引理~\ref{lem:tensor-integrality}、
递归恒等式与恢复引理从叶节点出发归纳证明正确性。

一个调用处的插值节点数、取值数和组数关于 $s$ 是多项式的，且与 $\mathcal L$ 无关。
由 \eqref{eq:tensor-partition-count}，一个活动深度 $d$ 调用的子节点数至多为
\[
 s^{O(1)}\exp\left(
       O_C\left(\frac{s}{\sqrt\lambda}
                    \left(\frac{\sqrt2}{8}\right)^d\right)\right).
\]
事实上 $M\le Cs\,8^{-d}$ 且
$u_d=\Theta_C(\lambda/2^d)$。将这些界在
$O_C(\log\lambda)$ 个深度上相乘，或以这种方式界定每一层上的节点数，得到总调用数至多为
\[
 \exp\left(O_C(s/\sqrt\lambda)
                   +O_C(\lambda\log\lambda)\right)
       =\exp(O_C(s/\sqrt\lambda))
\]
次。在叶节点处，直接求和至多使用 $2^H$ 个赋值，且每个赋值的项计算次数为多项式。
其代价被 $s+\mathcal L$ 的多项式乘以
\[
 \exp(O_C(H)+O_{C,A}(s^{9/10}))
   =\exp(O_{C,A}(s/\sqrt\lambda)).
\]
所界定。局部算术、打包、描述和恢复，除已界定的项计算之外，每生成一项
只使用多项式时间。依赖于描述的因子界定每次调用的工作量。由于分支界与
$\mathcal L$ 无关，总工作量保持 $s+\mathcal L$ 的固定多项式因子。
这证明了时间界。

关于空间，深度优先遍历递归，每次累加一个关于 $\beta$ 的和与一个关于组的乘积。
存储一个调用恢复所需的多项式多个 Taylor 数据、其当前的插值与模式指标，以及其
图与新张量描述。它们的大小关于 $s+\mathcal L$ 是多项式的，其中精度由
\eqref{eq:tensor-precision-bound} 界定。栈的层数只有
$O_C(\log\lambda)$。引理~\ref{lem:tapered-hermite} 中的公分母
计算也只使用多项式位整数。每个调用在其自身固定精度下累加其收缩项。
原始项计算例程按假设使用多项式空间，
新例程具有同样的性质。
这证明了空间界。最后在 $\mathcal A_{P_0}$ 中提取基
坐标，它们给出 $\mathcal B_{P_0,\ell}$ 中的答案。
\end{proof}
